\begin{abstract}
Network-traffic measurements describe communication patterns, but similar measurements do not necessarily identify the same attack category. The study examined highly imbalanced traffic data containing 10 classes, using a custom batch self-organising map trained without target labels. Clustering grouped the stored representative vectors before training labels mapped the clusters to classes. The frozen confirmatory classifier achieved assessment macro $F_1$ of \ConfirmF\ on folds that kept identical feature vectors together, while geometric diagnostics remained comparable between training and held-out observations. Full class-frequency correction increased rare-class recall but produced thousands of false positives. After the original assessment had been inspected, post-hoc exploration combined finer clustering with partial correction, reaching macro $F_1$ of \ExploreBF\ at 400 clusters. The exploratory gain accompanied greater codebook fragmentation and lower cluster compactness without providing independent confirmation. Nearest-neighbour and classification-tree baselines retained higher macro $F_1$ of \NeighbourF\ and \TreeF\ on identical folds. The findings showed that useful unsupervised geometry did not ensure reliable attack prediction, while class-frequency correction introduced a substantial precision--recall trade-off.
\end{abstract}
